\documentclass[11pt]{article}
\usepackage[T1]{fontenc}
\usepackage{lmodern,amsmath,amssymb,amsthm,booktabs}
\usepackage[margin=1in]{geometry}
\usepackage[colorlinks=true,linkcolor=blue,citecolor=blue,urlcolor=blue]{hyperref}
\newtheorem{theorem}{Theorem}
\newtheorem{lemma}{Lemma}
\newtheorem{proposition}{Proposition}
\newcommand{\norm}[1]{\left\lVert#1\right\rVert}
\newcommand{\nnz}{\operatorname{nnz}}
\newcommand{\round}{\operatorname{round}}
\newcommand{\softO}{\widetilde O}
\title{A rational-input specialization of Kelner's sparse linear solver}
\author{Kaiwen Jin\\\small\href{mailto:yui.ui.siki@gmail.com}{yui.ui.siki@gmail.com}}
\date{12 October 2026}
\hypersetup{pdftitle={A rational-input specialization of Kelner's sparse linear solver},pdfauthor={Kaiwen Jin}}
\begin{document}
\maketitle
\begin{abstract}
Kelner's recent identity-modulus lifting method gives a deterministic route to
the quadratic bit-work scale for sparse linear systems. We specialize this
method to the short rational inputs in AIM problem 044. For each fixed positive
integer $c$, a nonsingular rational $n\times n$ matrix with $m$ nonzeros,
condition number at most $n^c$, and input numerators and denominators of
$O_c(\log(2n))$ bits can be solved to relative Euclidean solution error at most
$n^{-c}/10$ in $\softO_c(nm)$ bit operations. Individual rational rounding
avoids a common denominator over the input. The proof gives explicit scales,
short residual bounds, reconstruction heights, and a direct forward-error
estimate. An accompanying implementation uses FLINT for fast bounded rational
reconstruction and includes exact finite-case checks.
\end{abstract}

\section{The promised-input problem}
The distinction between arithmetic operations and bit operations matters for
sparse linear systems: even an algorithm with few matrix--vector products may
operate on increasingly long numbers. The finite-precision question appears
as Problem 2.10 of the Simons workshop collection~\cite{workshop}. AIM
problem 044~\cite{aim} fixes a polynomial condition bound and short rational
input, and asks for $\softO(nm)$ bit operations at inverse-polynomial relative
solution accuracy.

Kelner~\cite{kelner} introduced a deterministic solver using a perturbation
whose scaled matrix is the identity modulo a chosen integer. Its lifting
steps need only sparse products on short integers; fast rational
reconstruction recovers the solution of the perturbed system. That method
and its fast reconstruction theorem are the basis of this note. Our purpose
is to make its consequence for the exact rational-input and forward-error
formulation of problem 044 explicit. The algorithmic mechanism is due to
Kelner.

All vector norms without a subscript are Euclidean, and matrix norms with
subscript $2$ are induced Euclidean norms. We use
$\kappa_2(A)=\norm{A}_2\norm{A^{-1}}_2$. A rational number is represented by
a signed numerator and a positive denominator. Bit work includes arithmetic,
input and output, and sparse indices. The notation $\softO_c$ hides fixed
powers of $\log(2n)$, with constants allowed to depend on the fixed integer
$c$.

\begin{theorem}\label{thm:main}
For every fixed positive integer $c$ there is a deterministic algorithm with
the following guarantee. Let $n\geq1$, let $A\in\mathbb Q^{n\times n}$ be
nonsingular and given by its sparse nonzero list, put $m=\nnz(A)$, and let
$b\in\mathbb Q^n\setminus\{0\}$. Suppose
\[
 \kappa_2(A)\leq n^c,
 \qquad B=c\lceil\log_2(2n)\rceil,
\]
and every input numerator and denominator has at most $B$ bits. In
$\softO_c(nm)$ bit operations the algorithm returns a rational vector
$\widehat x$ satisfying
\begin{equation}\label{eq:target}
 \norm{\widehat x-A^{-1}b}_2
 \leq \frac{n^{-c}}{10}\norm{A^{-1}b}_2.
\end{equation}
\end{theorem}
The probability of success is one, which meets the randomized success
requirement of problem 044. Symmetry, positive definiteness, and graph
structure are not required. Nonsingularity and the condition bound are input
promises; the algorithm need not certify them.

\section{An exact integer solve by lifting}
We first describe the integer part of the construction. This is Kelner's
identity-modulus method, with convenient explicit bounds. For an integer
matrix $C\in\mathbb Z^{n\times n}$ and $d\in\mathbb Z^n$, set
\[
 T=\max\{1,\max_{i,k}|C_{ik}|,\max_i|d_i|\}.
\]
Choose a power of two $R=2^s\geq2nT$ and put $M=I+RC$. The determinant
expansion gives $\det M\equiv1\pmod R$. Thus $M$ is nonsingular, and its
determinant is coprime to every power of $R$. In particular this construction
does not require $C$ to be nonsingular.

To solve $Mz=d$, fix $L=6n+2$, start with $r_0=d$, and compute, for
$j=0,\ldots,L-1$,
\begin{equation}\label{eq:lift}
 u_j=r_j\bmod R\in\{0,\ldots,R-1\}^n,
 \qquad
 r_{j+1}=\frac{r_j-u_j}{R}-Cu_j.
\end{equation}
Remainders are taken coordinatewise and are nonnegative even for negative
residuals. The divisions in~\eqref{eq:lift} are exact. The digits encode the
modular solution; they are stored during the iteration and assembled only
after the last step.

\begin{lemma}\label{lem:lift}
Let $P=R^L$ and $v=\sum_{j=0}^{L-1}R^ju_j$. Then $0\leq v_i<P$ and
$Mv\equiv d\pmod P$. Each residual coordinate has absolute value at most
$R^2$, and all arithmetic operands used in the sparse products have $O(s)$
bits. The lifting and assembly cost
$O(n(\nnz(C)+n)\operatorname{poly}(s+\log(2n)))$ bit operations.
\end{lemma}
\begin{proof}
Equation~\eqref{eq:lift} gives $Mu_j=r_j-Rr_{j+1}$. Multiplication by $R^j$
and summation telescope to
\begin{equation}\label{eq:telescoping}
 Mv=d-R^Lr_L.
\end{equation}
The digit bounds give the asserted range of $v$.

For a vector, $\norm{\cdot}_\infty$ denotes its largest absolute coordinate;
for a matrix it denotes the largest absolute row sum. Put
$W=\max\{\norm{d}_\infty,\norm{M}_\infty\}$. If
$\norm{r_j}_\infty\leq W$, then
\[
 \norm{r_{j+1}}_\infty
 \leq\frac{W+(R-1)\norm{M}_\infty}{R}\leq W.
\]
This starts at $r_0=d$. Since $nT\leq R/2$,
\[
 \norm{M}_\infty\leq1+nRT\leq1+R^2/2\leq R^2,
 \qquad \norm{d}_\infty\leq T\leq R^2.
\]
Every partial sum in a row of $Cu_j$ is bounded in absolute value by
$nT(R-1)\leq R^2/2$. Residuals, digits, products, sums, and the exact
division therefore involve $O(s)$ bits. Each of the $O(n)$ rounds costs
$O(\nnz(C)+n)$ short-integer operations, with the sparse indices charged.
Elementary arithmetic suffices at these widths. As $R$ is a power of two,
concatenating the $s$-bit digits once per coordinate constructs $v$ in
$O(nLs)$ bit operations. Repeatedly updating long prefixes is unnecessary.
\end{proof}

The last step uses bounded rational reconstruction. We state the precise
external result because its arithmetic cost is essential.

\begin{proposition}[Kelner, Theorem 2.3]\label{prop:recovery}
Let $P,v,N,D$ be integers with $P\geq2$, $0\leq v<P$, $N,D\geq1$, and
$P>2ND$. There is at most one reduced fraction $a/q$ satisfying
\[
 |a|\leq N,\quad 1\leq q\leq D,\quad
 \gcd(q,P)=1,\quad a\equiv vq\pmod P.
\]
If it exists, it can be recovered deterministically in
$O(\mathsf M(k)\log k)$ bit operations, where
$k=\lceil\log_2(P+1)\rceil$ and $\mathsf M(k)=\softO(k)$ is the bit cost of
fast integer multiplication.
\end{proposition}
This is Theorem 2.3 on page 7 of the pinned version~\cite{kelner}; its modulus
need not be prime. The fast reconstruction algorithm is part of this
statement. Ordinary scalar Euclidean reconstruction cannot replace it in
the bit-work analysis.

\begin{lemma}\label{lem:height}
With $M,d,R,L,P,v$ as above, set $H=R^{3n}$. Applying
Proposition~\ref{prop:recovery} to each $v_i$ with $N=D=H$ recovers exactly
$z=M^{-1}d$. All coordinates together cost $\softO(n^2s)$ bit operations,
including parameter construction and output writing.
\end{lemma}
\begin{proof}
Write $z_i=a_i/q_i$ in lowest terms, with $q_i>0$. By Cramer's rule,
$z_i=\det(M_i)/\det M$, where $M_i$ replaces column $i$ by $d$. Hence
$q_i$ divides $|\det M|$ and is coprime to $P$.

All entries in $M$ and $d$ have absolute value at most $RT+1\leq R^2$.
Hadamard's determinant bound gives
\[
 |\det M|,\ |\det(M_i)|
 \leq n^{n/2}R^{2n}\leq R^{5n/2}\leq H,
\]
where $n\leq R$. For completeness, orthogonalizing the columns preserves
the determinant and does not increase their lengths, so its absolute value
is bounded by the product of the original column lengths. A dependent
column gives determinant zero. Reduction of the Cramer fractions therefore
preserves $|a_i|\leq H$ and $q_i\leq H$.

Since $M$ is invertible modulo $P$,~\eqref{eq:telescoping} implies
$a_i\equiv v_iq_i\pmod P$. Also
\[
 P=R^{6n+2}\geq4R^{6n}>2H^2.
\]
All hypotheses of Proposition~\ref{prop:recovery} are now satisfied,
including existence and coprimality. The modulus has $O(ns)$ bits, giving
$\softO(ns)$ work per coordinate. The parameters $P,H$ are binary powers
of two, constructed by writing their bits. Every output numerator and
denominator has $O(ns)$ bits, and their total writing cost is $O(n^2s)$.
\end{proof}

\section{Rational rounding and forward error}
We now choose the integer data for Theorem~\ref{thm:main}. Put
$Q=2^B$ and $K=n^c$. Each nonzero input has absolute value at least $1/Q$
and each input has absolute value at most $Q$. Nonsingularity gives a
nonzero matrix entry, and a nonzero column tested on a coordinate vector
therefore shows $\norm{A}_2\geq1/Q$. Consequently
\begin{equation}\label{eq:inputscale}
 \norm{A^{-1}}_2\leq KQ,\qquad
 \norm{b}_2\geq1/Q,
 \qquad x=A^{-1}b\ne0.
\end{equation}
Every row of $A$ is nonzero, so $m\geq n$.

To keep the perturbation small even at the largest allowed condition
number, choose
\begin{equation}\label{eq:roundscale}
 \ell=\lceil\log_2 n\rceil,\quad
 h=4+B+(2c+1)\ell,\quad S=2^h,
 \qquad S\geq16K^2nQ.
\end{equation}
Round entries individually by
\[
 C_{ik}=\round(SA_{ik}),\qquad d_i=\round(Sb_i),
 \qquad \round(t)=\lfloor t+1/2\rfloor.
\]
On an input $p/q$ with $q>0$, the rounded value is
$\lfloor(2Sp+q)/(2q)\rfloor$. This requires only $O_c(\log(2n))$-bit
operands. Zeros remain zero, so $\nnz(C)\leq m$, and no common input
denominator is formed. Entrywise rounding errors give
\begin{equation}\label{eq:rounderror}
 \norm{C/S-A}_2\leq\frac{n}{2S},\qquad
 \norm{d/S-b}_2\leq\frac{\sqrt n}{2S}.
\end{equation}
The first inequality follows by bounding the operator norm by the
Frobenius norm, using Cauchy--Schwarz on each row.

For these $C,d$, the quantity $T$ in Section 2 obeys $T\leq SQ+1/2$, so
$\log_2 T=O_c(\log(2n))$. Choose the least power of two $R\geq2nT$, form
$M=I+RC$, and use the lifting and reconstruction of Section 2 to obtain
$z=M^{-1}d$. Return
\begin{equation}\label{eq:output}
 \widehat x=Rz.
\end{equation}
The complete computation is summarized in the algorithm below.

\begin{center}
\fbox{\begin{minipage}{0.94\linewidth}
\textbf{Algorithm: sparse rational solve}\par\smallskip
\textbf{Input:} $n,c,A,b$ satisfying Theorem~\ref{thm:main}.\par
\textbf{Output:} $\widehat x\in\mathbb Q^n$.\par\smallskip
\begin{tabbing}
\hspace{1em}\=\hspace{1em}\=\kill
Compute $B,\ell,h,S$ by~\eqref{eq:roundscale} and the definition of $B$.\\
Set $C=\round(SA)$ and $d=\round(Sb)$ entrywise.\\
Set $T=\max\{1,\max|C_{ik}|,\max|d_i|\}$; choose $R=2^{\lceil\log_2(2nT)\rceil}$.\\
Set $L=6n+2$ and $r=d$.\\
\textbf{for} $j=0,\ldots,L-1$ \textbf{do}\\
\> $u_j=r\bmod R$; store the $n$ radix-$R$ digits.\\
\> $w=Cu_j$ using the sparse nonzero list of $C$.\\
\> $r=(r-u_j)/R-w$.\\
\textbf{end for}\\
Set $P=R^L$ and $H=R^{3n}$ as binary powers of two.\\
\textbf{for} $i=1,\ldots,n$ \textbf{do}\\
\> Concatenate the stored digits to form $v_i=\sum_{j=0}^{L-1}R^j(u_j)_i$.\\
\> Recover $a_i/q_i$ from $v_i,P,N=D=H$ by Proposition~\ref{prop:recovery}.\\
\> $\widehat x_i=Ra_i/q_i$.\\
\textbf{end for}\\
\textbf{return} $\widehat x$.
\end{tabbing}
\end{minipage}}
\end{center}

Lemmas~\ref{lem:lift} and~\ref{lem:height} give
$\softO_c(n(m+n)+n^2)=\softO_c(nm)$ bit work for this computation, including
rounding, assembly, reconstruction and output. Stored digits and sparse data,
with fast reconstruction workspace, require $\softO_c(n^2)$ bits.
Diagnostics and exact reference solutions used in the tests are outside the
solver. The following proof establishes its relative error and collects the
work estimates.

\begin{proof}[Proof of Theorem~\ref{thm:main}]
The exact working equation $Mz=d$ becomes
\[
 (A+E)\widehat x=b+f,
 \qquad E=C/S-A+I/(RS),\quad f=d/S-b.
\]
Both perturbations must be included in the forward error. By
\eqref{eq:rounderror} and $R\geq2$,
\[
 \norm{E}_2\leq n/S,\qquad \norm{f}_2\leq n/(2S).
\]
Let $\beta=\norm{A^{-1}}_2\norm{E}_2$. Equations
\eqref{eq:inputscale}--\eqref{eq:roundscale} yield
\[
 \beta\leq\frac{KnQ}{S}\leq\frac1{16K}\leq\frac1{16}.
\]
For any vector $w$, the triangle inequality gives
$\norm{(I+A^{-1}E)w}_2\geq(1-\beta)\norm{w}_2$. This square matrix is
injective, hence invertible, with inverse norm at most $(1-\beta)^{-1}$.
Subtracting the equation $Ax=b$, and writing $e=\widehat x-x$, gives
\[
 (I+A^{-1}E)e=A^{-1}f-A^{-1}Ex.
\]
As $\norm{b}_2=\norm{Ax}_2\leq\norm{A}_2\norm{x}_2$,
\[
 \frac{\norm{A^{-1}f}_2}{\norm{x}_2}
 \leq\kappa_2(A)\frac{\norm{f}_2}{\norm{b}_2}
 \leq\frac{KnQ}{2S}.
\]
Combining these estimates proves the explicit bound
\begin{equation}\label{eq:forward}
 \frac{\norm{e}_2}{\norm{x}_2}
 \leq\frac{3KnQ/(2S)}{1-KnQ/S}
 \leq\frac{3}{2(16K-1)}\leq\frac1{10K}.
\end{equation}
The last inequality is equivalent to $30K\leq32K-2$, valid for $K\geq1$.
This proves~\eqref{eq:target}, including $n=1$.

For the bit work, $h$ and $s=\log_2 R$ are $O_c(\log(2n))$. Rounding and
reading the input cost $O_c((m+n)\operatorname{polylog}(2n))$.
Lemma~\ref{lem:lift} gives $\softO_c(n(m+n))$ lifting and digit-assembly
work. Lemma~\ref{lem:height} gives $\softO_c(n^2)$ reconstruction work.
Multiplying the recovered numerators by $R$ and writing the rational output
also cost $O(n^2s)$. The recovered denominators are odd, so this multiplication
preserves reduced fractions without a further long-integer gcd. Since
$m\geq n$, the total cost is $\softO_c(nm)$. Every operation is deterministic.
\end{proof}

\section{Implementation and finite checks}
The accompanying Python library accepts a sparse list with zero-based
indices and rational entries. Its default reconstruction backend calls
FLINT's bounded rational reconstruction, through the pinned
\texttt{python-flint 0.8.0} wheel~\cite{flint}. The binding supports its
64-bit Linux and macOS ABI. Integer transfer uses hexadecimal strings,
avoiding a decimal-conversion cost on long integers. Digits are stored in
byte arrays and concatenated once. For this purpose the code rounds $s$
up to a multiple of eight; the resulting $R$ is at most $2^7$ times the
least admissible radix. Every bound above remains valid because
$R\geq2nT$ and $s=O_c(\log(2n))$ still hold. A separate classical Euclidean
backend serves as a slow reference and carries no fast bit-complexity claim.

The tests compare the output with exact rational Gaussian elimination,
exercise negative and zero coordinates, non-dyadic denominators, small input
scales, and the condition-bound endpoint, and compare fast reconstruction
with the Euclidean reference. A fixed 32,768-bit reconstruction test checks
the long-integer interface. The reproduction script uses $c=2$ and
$n\in\{4,8,16,32\}$ in each of four families: a diagonal matrix with entries
$1,\ldots,1,n^{-2}$; a scalar matrix $I/(2n^2-1)$; a tridiagonal matrix with
alternating diagonal entries $\pm2$ and off-diagonal entries $1/4$; and an
upper bidiagonal matrix with diagonal $1$ and superdiagonal $1/4$.
The right-hand sides have coordinates $(-1)^i/(2i+1)$, indexed from zero.
The first two families have condition numbers $n^2$ and $1$. For the third,
the off-diagonal operator has norm at most $1/2$, giving singular values in
$[3/2,5/2]$; for the fourth they lie in $[3/4,5/4]$. Thus all cases meet the
input promises.

All 14 tests and all 16 reproduced cases passed. Case acceptance uses the
exact squared form of~\eqref{eq:target}, with an independent rational
solution as reference, together with the proved residual-width bound.
Code, compact results, and reproduction instructions are available in the
\href{https://github.com/yuiyuiui/aim-044-sparse-lift/tree/d1541aa86177584bcf75dc8545d8859a4d549258}{fixed version of the accompanying repository}. These
finite checks validate the implementation on the stated cases; the
all-input guarantee and asymptotic work bound are established by the proof.

\paragraph{AI use.}
OpenAI Codex assisted with literature retrieval, derivation, proof checking,
implementation, finite-case validation, and preparation of this manuscript.
The central algorithm and the fast rational reconstruction result are
attributed to Kelner. This statement does not assert independent human review
or formal verification.

\begin{thebibliography}{9}
\bibitem{aim}
M.~Colbrook and contributors,
\emph{044. Quadratic-scale bit complexity for well-conditioned sparse systems},
AIM problem catalogue, revision
\texttt{c9929805cece044705abaaebe1cebddac1ba11a6}, 2026.
\href{https://github.com/MColbrook/AIM/blob/c9929805cece044705abaaebe1cebddac1ba11a6/problems/044-sparse-linear-solve-bit-complexity.md}{Pinned problem statement}.
\bibitem{workshop}
N.~Amsel et al.,
\emph{Linear Systems and Eigenvalue Problems: Open Questions from a Simons Workshop},
arXiv:2602.05394, 2026, Problem 2.10.
\url{https://arxiv.org/abs/2602.05394}.
\bibitem{kelner}
J.~A.~Kelner,
\emph{Solving Linear Systems in $\softO(mn\log(\kappa/\epsilon))$ Bit Operations},
arXiv:2609.38101v1, 29 September 2026.
\url{https://arxiv.org/abs/2609.38101v1}.
\bibitem{flint}
FLINT contributors,
\emph{FLINT 3.3.1: Fast Library for Number Theory},
bounded rational reconstruction routine
\texttt{fmpq\_reconstruct\_fmpz\_2}, 2025.
\href{https://github.com/flintlib/flint/blob/v3.3.1/src/fmpq/reconstruct_fmpz_2.c}{Versioned source};
\href{https://pypi.org/project/python-flint/0.8.0/}{python-flint 0.8.0}.
\end{thebibliography}
\end{document}
